% =====================================================================
\chapter{Estado del arte}\label{cap:estado-arte}
% =====================================================================

El presente trabajo se sitúa en la intersección de tres líneas de investigación activas: los efectos de los edulcorantes no calóricos sobre el neurodesarrollo, el impacto del consumo temprano de azúcar sobre la neurogénesis hipocampal, y la aplicación de métodos de aprendizaje automático al fenotipado conductual. La revisión de la literatura reciente revela un panorama caracterizado por resultados heterogéneos y, en el caso particular de la stevia, abiertamente contradictorios, lo que fundamenta la pertinencia de una caracterización sistemática.

% ---------------------------------------------------------------------
\section{Edulcorantes no calóricos y programación del neurodesarrollo}\label{sec:edulcorantes}

El consumo de edulcorantes no calóricos se ha extendido a poblaciones sensibles, incluidas embarazadas y niños, y la evidencia sugiere que no son metabólicamente inertes. Estudios de neuroimagen muestran que estos compuestos modifican la respuesta de regiones cerebrales vinculadas a la recompensa y la saciedad, y que la exposición durante el desarrollo temprano puede alterar las preferencias por el sabor dulce \cite{yunker2020nns}. En modelos animales, los efectos trascienden al individuo expuesto: la exposición a sacarina en ratones macho produjo impulsividad motora, hiperactividad locomotora y déficits de memoria de trabajo en la descendencia, asociados a hipermetilación del ADN espermático en regiones promotoras de receptores dopaminérgicos \cite{mccarthy2020saccharin}. En la misma línea, la ingesta materna elevada de endulzantes ---incluida la stevia--- durante la gestación y la lactancia modificó la microbiota intestinal y agravó los procesos de aprendizaje y memoria en las crías macho adultas \cite{delagarza2022maternal}. Trabajos con exposición juvenil directa muestran que el consumo temprano de edulcorantes de bajas calorías altera la regulación de la glucosa, la conducta motivada por el azúcar y la memoria episódica dependiente del hipocampo, con cambios en vías de señalización génica hipocampal de manera dependiente del sexo \cite{tsan2022lcs}.

% ---------------------------------------------------------------------
\section{\textit{Stevia rebaudiana}: evidencia contradictoria sobre el sistema nervioso}\label{sec:stevia-evidencia}

La stevia ocupa una posición ambigua en la literatura. Por un lado, se le atribuyen propiedades neuroprotectoras: en ratas alimentadas con una dieta alta en grasas y sacarosa, la suplementación con \textit{Stevia rebaudiana} redujo los marcadores de peroxidación lipídica cerebral en un 30--51\,\% \cite{pozdnyakova2025stevia}, y los glucósidos de esteviol han demostrado efectos insulino-miméticos y antioxidantes, incrementando las defensas enzimáticas frente al estrés oxidativo \cite{prata2017glycosides}. En modelos patológicos, el esteviósido mostró un efecto antiamnésico en ratas tratadas con escopolamina \cite{singh2010stevioside}, la stevia modificó la plasticidad sináptica y los niveles de NADPH oxidasa del sistema nervioso central en un modelo de desorden metabólico inducido por fructosa \cite{chavushyan2017stevia}, y su administración revirtió la hiperglucemia, la ansiedad y el deterioro de memoria en ratas diabéticas \cite{khakpai2023stevia}. Por otro lado, la evidencia conductual es menos favorable: un estudio en ratas evaluó el consumo de edulcorantes de bajas calorías ---entre ellos la propia stevia--- durante la adolescencia y halló deterioros persistentes en la memoria dependiente del hipocampo en la adultez \cite{tsan2022lcs}. Los trabajos del laboratorio anfitrión se inscriben en esta tensión: la stevia afecta la función hipocampal y activa el sistema dopaminérgico de forma análoga a otras recompensas naturales, aunque algunos autores cuestionan estos hallazgos por las dosis empleadas o por la imposibilidad de replicarlos con esteviósidos purificados \cite{kruse2025stevia}. En ratas, el consumo de stevia se asoció además con una menor expresión de $\Delta$FosB, un marcador de activación neuronal crónica, en el sistema dopaminérgico \cite{salinasvelarde2021fosb}. Esta discrepancia ---agravada por la ausencia de estudios en ratas macho con una batería conductual e histológica combinada en las ventanas etarias juvenil y adulta--- constituye el principal vacío que este proyecto aborda.

% ---------------------------------------------------------------------
\section{Consumo temprano de azúcar, neurogénesis hipocampal y separación de patrones}\label{sec:azucar-neurogenesis}

La línea de la sacarosa aporta el marco de referencia contra el cual contrastar los efectos de la stevia. El consumo de dietas ricas en azúcar durante las etapas juvenil y adolescente deteriora selectivamente la función hipocampal, preservando en algunos casos formas de memoria no dependientes del hipocampo \cite{reichelt2016sucrose,coirini2022sucrose}. El acceso temprano a sacarosa tuvo efectos negativos de largo plazo sobre la memoria en ratas macho \cite{noble2019early}, y el consumo de sacarosa o de jarabe de maíz de alta fructosa durante la adolescencia se estudió en relación con la memoria espacial y la neuroinflamación hipocampal \cite{hsu2015sucrose}, con el control cognitivo, la memoria de reconocimiento y la inmunorreactividad de parvalbúmina \cite{reichelt2015adolescent}, y con la memoria, la conducta tipo ansiosa y los niveles hipocampales de parvalbúmina y doblecortina \cite{xu2018sucrose}. A nivel mecanístico, se ha demostrado que las alteraciones de la microbiota inducidas por el azúcar son causalmente responsables de los déficits de memoria, mediante el enriquecimiento de taxones específicos del género \textit{Parabacteroides} \cite{noble2021microbial}, y que el consumo crónico de sacarosa reduce en más del 40\,\% la expresión de genes hipocampales vinculados a la plasticidad y la neurogénesis, con efectos comparables a los del estrés temprano \cite{maniam2016sugar}. No obstante, los resultados dependen críticamente de la ventana de exposición: el consumo de sacarosa restringido a la adolescencia tardía no indujo conducta ansiosa y, contra lo esperado, aumentó la proliferación celular en la zona subgranular, si bien con neuronas inmaduras de orientación dendrítica alterada \cite{sanchezhuerta2022sucrose}. Esta sensibilidad al momento de exposición respalda directamente la elección de dos ventanas etarias contrastantes (PD25 y PD75) en el diseño experimental.

% ---------------------------------------------------------------------
\section{Análisis multivariado y aprendizaje automático en el fenotipado conductual}\label{sec:ml-fenotipado}

El análisis conductual tradicional se apoya en medidas discretas y univariadas que, examinadas de forma aislada, omiten la estructura de orden superior contenida en el comportamiento animal \cite{geros2020tracking}. Frente a esta limitación, el aprendizaje automático se ha consolidado como herramienta de fenotipado: se ha demostrado que algoritmos como Random Forest, máquinas de soporte vectorial y regresión logística clasifican estados conductuales de roedores con exactitudes del 85--96\,\% \cite{gharagozloo2021ml}. En aplicaciones específicas, los modelos de Random Forest han permitido discriminar fenotipos asociados a disfunción dopaminérgica de aquellos de origen serotoninérgico, e identificar las variables conductuales de mayor peso discriminativo \cite{lukovikov2026cylinder}, y clasificadores más simples han diferenciado cepas de ratas a partir de sus perfiles en el campo abierto y el laberinto elevado \cite{geros2020tracking}. Sin embargo, la aplicación de estas técnicas a datos conductuales provenientes de intervenciones nutricionales permanece escasamente explorada, lo que sitúa el componente de PCA, clustering jerárquico y Random Forest de este proyecto como un aporte metodológico diferencial desde la bioingeniería.
